\chapter{Nombre dérivé, tangente et approximation}
\section*{Exercice 8.1 -- Revenir à la définition}
Pour $h\ne0$, $[(-2+h)^2-4]/h=(-4h+h^2)/h=-4+h$. La limite quand $h\to0$ vaut $-4$, donc $f'(-2)=-4$.

\section*{Exercice 8.2 -- Tangente}
$f(2)=1/2$ et $f'(2)=-1/4$. L'équation est $y=1/2-(x-2)/4=1-x/4$. Pour $x=2$, on retrouve $y=1/2$. Le nombre dérivé est une pente, pas l'ordonnée à l'origine.

\section*{Exercice 8.3 -- Racine à l'origine}
Le taux vaut $\sqrt h/h=1/\sqrt h$. Il tend vers $+\infty$ lorsque $h\to0^+$. Une dérivée réelle exige une limite finie. La courbe peut avoir une tangente verticale sans posséder un nombre dérivé fini.

\section*{Exercice 8.4 -- Valeur absolue}
Le taux $|h|/h$ vaut $1$ pour $h>0$ et $-1$ pour $h<0$. Les deux limites unilatérales diffèrent ; la dérivée n'existe pas. La fonction est pourtant continue, puisque $|h|\to0$.

\section*{Exercice 8.5 -- Approximation affine}
La tangente donne $2+0{,}04/4=2{,}01$. L'identité exacte $\sqrt{4+h}-2-h/4=-h^2/[4(\sqrt{4+h}+2)^2]$ montre que l'approximation est trop grande. Pour $h=0{,}04$, l'erreur absolue est au plus $h^2/64=0{,}000025$.

\section*{Exercice 8.6 -- Vitesse et unités}
$s(2+h)-s(2)=7h+h^2$. La vitesse moyenne est $7+h$ mètres par seconde pour $h\ne0$ admissible. Sa limite vaut $7$ m/s. La position $s(2)=10$ m et la vitesse $7$ m/s sont des quantités différentes.
